% preamble.tex —— 共享导言区（不含 \documentclass，由 gen_sections.py 内联进 main.tex）
% 注意：不引入 ocgx2 或任何 OCG 图层包；hyperref 及导航宏在下方注入。

\usepackage{amsmath,amssymb,amsthm,mathtools}
\usepackage{amsfonts}
\usepackage{esint}   % 提供 \oiint、\oiiint 等闭曲面积分符号（与 conventions 符号表对齐）
\usepackage{bm}        % \bm 粗体数学（向量/张量），DeepSeek 常用
\usepackage{mathrsfs}  % \mathscr 花体（傅里叶变换、拉氏变换等）
\usepackage[a4paper,margin=2.4cm]{geometry}
\usepackage{enumitem}
\usepackage{xcolor}
\usepackage{array}
% siunitx（单位/科学计数，大学物理常用）按需加载——缺包时跳过而非整卷编译失败
\IfFileExists{siunitx.sty}{\usepackage{siunitx}}{}
% TikZ / PGFPlots（矢量配图，--figures 时启用）——同样用 \IfFileExists 缺包兜底。
% 二者均不需要 shell-escape，与 -no-shell-escape 微编译及整卷编译完全兼容。
\IfFileExists{tikz.sty}{%
  \usepackage{tikz}%
  % decorations.markings：场线/有向路径中点箭头（物理电磁场常用）；bending：箭头随路径弯曲；
  % backgrounds：背景层；shapes.geometric/shapes.misc：几何/圆角节点；fit：包裹多节点的框。
  \usetikzlibrary{calc,arrows.meta,angles,quotes,patterns,intersections,positioning,decorations.pathmorphing,decorations.markings,bending,backgrounds,shapes.geometric,shapes.misc,fit}%
}{}
% pgfplots 库必须在 \usepackage{pgfplots} 之后加载。fillbetween：两曲线间填充（阴影概率/积分区域，
% 各科政策都会用到）；groupplots：多子图；polar：极坐标；statistics：箱线/直方等统计图。
% 缺这些库时，模型一旦按政策用 fill between 等就会整卷编译失败——这正是配图不稳定的主因。
\IfFileExists{pgfplots.sty}{%
  \usepackage{pgfplots}\pgfplotsset{compat=1.18}%
  \usepgfplotslibrary{fillbetween,groupplots,polar,statistics}%
}{}

% ---- 防御性兜底：补常见但默认未定义的宏，避免 "Undefined control sequence" 整卷失败 ----
% 一律用 \providecommand（仅在未定义时生效），不覆盖任何包或各科 conventions 的既有记号。
\providecommand{\R}{\mathbb{R}}
\providecommand{\N}{\mathbb{N}}
\providecommand{\Z}{\mathbb{Z}}
\providecommand{\Q}{\mathbb{Q}}
\providecommand{\C}{\mathbb{C}}
\providecommand{\dd}{\,\mathrm{d}}   % 微分 d
\providecommand{\degree}{^\circ}
\providecommand{\abs}[1]{\left|#1\right|}
\providecommand{\norm}[1]{\left\|#1\right\|}
% \sloppy：放宽断行，减少 overfull（仅警告非错误，但让长公式版面更稳）
\sloppy

<<<HYPERLINKS_PREAMBLE>>>

\linespread{1.15}
\setlength{\parindent}{0pt}
\setlength{\parskip}{4pt}

% 最终答案框
\newcommand{\boxedans}[1]{\[\boxed{\,#1\,}\]}

% 常见陷阱提示行
\newcommand{\pitfall}[1]{\par\smallskip\noindent\textcolor{red!60!black}{\textbf{常见陷阱：}}#1\par}

% 知识点小节标题（study 模式用）
\newcommand{\kpsection}[1]{\par\bigskip\noindent{\large\bfseries #1}\par\smallskip}
